% ==============================================================================
\section{Псевдообратная матрица Мура--Пенроуза}
\label{sec:pseudo-inverse}
% ==============================================================================

Для прямоугольной или вырожденной матрицы $\bm{A} \in \C^{m \times n}$ классическая обратная матрица не существует, а система линейных алгебраических уравнений $\bm{A}\bm{x} = \bm{b}$ может не иметь точного решения. В таком случае ставится задача поиска вектора $\bm{x}$, минимизирующего невязку:
\begin{equation}
    \min_{\bm{x} \in \C^n} \norm{\bm{A}\bm{x} - \bm{b}}_2.
    \label{eq:least-squares}
\end{equation}
Множество решений задачи~\cref{eq:least-squares} имеет вид $\bm{x} = \bm{A}^+ \bm{b} + (\bm{I} - \bm{A}^+ \bm{A})\bm{y}$, где $\bm{y} \in \C^n$ — произвольный вектор, а матрица $\bm{A}^+$ называется псевдообратной. Слагаемое $(\bm{I} - \bm{A}^+ \bm{A})$ представляет собой ортопроектор на ядро матрицы $\bm{A}$.

\begin{tcbdefinition}[Псевдообратная матрица через SVD]
\label{def:pinv-svd}
    Пусть $\bm{A} \in \C^{m \times n}$ — произвольная матрица с сингулярным разложением $\bm{A} = \bm{U}\bm{\Sigma}\bm{V}\H$, где $\bm{U} \in \C^{m \times m}$, $\bm{V} \in \C^{n \times n}$ — унитарные матрицы, а $\bm{\Sigma} \in \R^{m \times n}$ — матрица сингулярных чисел. \textit{Псевдообратной матрицей Мура--Пенроуза} называется матрица $\bm{A}^+ \in \C^{n \times m}$, определяемая как
    \begin{equation}
        \bm{A}^+ = \bm{V}\bm{\Sigma}^+\bm{U}\H,
        \label{eq:pinv-svd}
    \end{equation}
    где $\bm{\Sigma}^+ \in \R^{n \times m}$ получена из $\bm{\Sigma}$ транспонированием и заменой каждого ненулевого сингулярного числа $\sigma_i$ на $\sigma_i^{-1}$.
\end{tcbdefinition}

\begin{tcbremark}[Историческая справка]
    Понятие псевдообратной матрицы было введено Элиакимом Гастингсом Муром в 1920 году. Независимо от него концепцию разработали Арне Бьерхаммар в 1951 году и Роджер Пенроуз в 1955 году. В случае полноранговой матрицы ($\rank(\bm{A}) = n \leq m$) справедливо выражение $\bm{A}^+ = (\bm{A}\H\bm{A})^{-1}\bm{A}\H$.
\end{tcbremark}

\begin{tcbexample}[Простейшие случаи]
\label{ex:pinv-simple}
    \begin{itemize}
        \item Скалярный случай. Пусть $\bm{A} = [\alpha] \in \C^{1 \times 1}$. Если $\alpha \neq 0$, то $\bm{A}^+ = [1/\alpha]$. Если $\alpha = 0$, то $\bm{A}^+ = [0]$.
        \item Вектор-столбец. Пусть $\bm{A} = \begin{pNiceMatrix} 1 & 1 & 1 \end{pNiceMatrix}\T$. Тогда $\bm{A}\H\bm{A} = 3$, и
        \[
            \bm{A}^+ = \begin{pNiceMatrix} 1/3 & 1/3 & 1/3 \end{pNiceMatrix}.
        \]
        Умножение $\bm{x}^* = \bm{A}^+ \bm{b}$ для вектора правой части $\bm{b}$ дает среднее арифметическое элементов вектора $\bm{b}$.
    \end{itemize}
\end{tcbexample}

\begin{tcbwarning}[Ложное свойство мультипликативности]
\label{warn:ab-pinv}
    В отличие от обычной обратной матрицы, для псевдообратной матрицы в общем случае равенство $(\bm{A}\bm{B})^+ = \bm{B}^+\bm{A}^+$ \textbf{не выполняется}.
    
    Рассмотрим контрпример:
    \[
        \bm{A} = \begin{pNiceMatrix} 1 & 1 \\ 0 & 0 \end{pNiceMatrix}, \quad
        \bm{B} = \begin{pNiceMatrix} 0 & 0 \\ 1 & 1 \end{pNiceMatrix}.
    \]
    Псевдообратные матрицы равны:
    \[
        \bm{A}^+ = \begin{pNiceMatrix} 1/2 & 0 \\ 1/2 & 0 \end{pNiceMatrix}, \quad
        \bm{B}^+ = \begin{pNiceMatrix} 0 & 1/2 \\ 0 & 1/2 \end{pNiceMatrix}.
    \]
    Тогда $\bm{A}\bm{B} = \bm{A} \implies (\bm{A}\bm{B})^+ = \bm{A}^+$. С другой стороны,
    $\bm{B}^+\bm{A}^+ = \frac{1}{4} \begin{pNiceMatrix} 1 & 0 \\ 1 & 0 \end{pNiceMatrix}$. Очевидно, что $(\bm{A}\bm{B})^+ \neq \bm{B}^+\bm{A}^+$.
\end{tcbwarning}

\begin{tcbtheorem}[Алгебраическое определение Пенроуза]
\label{thm:penrose-conditions}
    Матрица $\bm{B} \in \C^{n \times m}$ является псевдообратной к $\bm{A} \in \C^{m \times n}$ тогда и только тогда, когда выполнены следующие четыре условия:
    \begin{enumerate}
        \item $\bm{A}\bm{B}\bm{A} = \bm{A}$.
        \item $\bm{B}\bm{A}\bm{B} = \bm{B}$.
        \item $(\bm{A}\bm{B})\H = \bm{A}\bm{B}$.
        \item $(\bm{B}\bm{A})\H = \bm{B}\bm{A}$.
    \end{enumerate}
    При этом для любой фиксированной матрицы $\bm{A}$ матрица $\bm{B}$ существует и единственна.
\end{tcbtheorem}

\begin{tcbproof}[теоремы~\ref{thm:penrose-conditions} (Единственность)]
    Предположим, существуют две матрицы $\bm{B}_1$ и $\bm{B}_2$, удовлетворяющие условиям 1--4. Докажем, что $\bm{B}_1 = \bm{B}_2$.
    \begin{align*}
        \bm{B}_1 &\overset{(2)}{=} \bm{B}_1 \bm{A} \bm{B}_1 \overset{(1)}{=} \bm{B}_1 \bm{A} \bm{B}_2 \bm{A} \bm{B}_1 \\
        &\overset{(3)}{=} \bm{B}_1 (\bm{A} \bm{B}_2)\H (\bm{A} \bm{B}_1)\H = \bm{B}_1 \bm{B}_2\H \bm{A}\H \bm{B}_1\H \bm{A}\H \\
        &= \bm{B}_1 \bm{B}_2\H (\bm{A} \bm{B}_1 \bm{A})\H \overset{(1)}{=} \bm{B}_1 \bm{B}_2\H \bm{A}\H \\
        &= \bm{B}_1 (\bm{A} \bm{B}_2)\H \overset{(3)}{=} \bm{B}_1 \bm{A} \bm{B}_2.
    \end{align*}
    Далее преобразуем полученное выражение, применяя свойства симметрично:
    \begin{align*}
        \bm{B}_1 \bm{A} \bm{B}_2 &\overset{(4)}{=} (\bm{B}_1 \bm{A})\H \bm{B}_2 = \bm{A}\H \bm{B}_1\H \bm{B}_2 \\
        &\overset{(1)}{=} (\bm{A} \bm{B}_2 \bm{A})\H \bm{B}_1\H \bm{B}_2 = \bm{A}\H \bm{B}_2\H \bm{A}\H \bm{B}_1\H \bm{B}_2 \\
        &= (\bm{B}_2 \bm{A})\H (\bm{B}_1 \bm{A})\H \bm{B}_2 \overset{(4)}{=} \bm{B}_2 \bm{A} \bm{B}_1 \bm{A} \bm{B}_2 \\
        &\overset{(1)}{=} \bm{B}_2 \bm{A} \bm{B}_2 \overset{(2)}{=} \bm{B}_2.
    \end{align*}
    Следовательно, $\bm{B}_1 = \bm{B}_2$.
\end{tcbproof}

% ==============================================================================
\section{Линейная регрессия и Total Least Squares}
\label{sec:tls}
% ==============================================================================

Пусть задан набор точек $(a_i, b_i) \in \R^2$. Задача классической линейной регрессии предполагает поиск коэффициента $x$ (и смещения $y$) для зависимости $a_i x + y \approx b_i$. В матричной форме задача формулируется как переопределенная система и сводится к минимизации:
\begin{equation}
    \min_{x, y} \norm{
    \begin{pNiceMatrix}
        a_1 & 1 \\
        \Vdots & \Vdots \\
        a_n & 1
    \end{pNiceMatrix}
    \begin{pNiceMatrix} x \\ y \end{pNiceMatrix}
    -
    \begin{pNiceMatrix} b_1 \\ \Vdots \\ b_n \end{pNiceMatrix}
    }_2^2.
    \label{eq:lin-reg}
\end{equation}
В~\cref{eq:lin-reg} минимизируется сумма квадратов \textit{вертикальных} отклонений от точек до прямой (вдоль оси ординат).

Альтернативный подход — \textit{Метод полных наименьших квадратов} (Total Least Squares, TLS). В нем минимизируется сумма квадратов кратчайших (перпендикулярных) расстояний от точек до прямой.

Пусть прямая задана нормальным вектором $\bm{v} = (v_1, v_2)\T$, где $\norm{\bm{v}}_2 = 1$. Кратчайшее расстояние от точки $(a_i, b_i)$ до такой прямой пропорционально скалярному произведению с нормалью. Задача TLS формализуется как поиск направления прямой, при котором проекции точек на нормаль минимальны:
\begin{equation}
    \min_{\norm{\bm{v}}_2 = 1} \sum_{i=1}^n (a_i v_1 + b_i v_2)^2 = \min_{\norm{\bm{v}}_2 = 1} \norm{\hat{\bm{A}} \bm{v}}_2^2,
    \label{eq:tls-opt}
\end{equation}
где $\hat{\bm{A}}$ составлена из координат точек. Данная задача решается через вычисление сингулярного разложения матрицы данных: оптимальный нормальный вектор $\bm{v}$ соответствует правому сингулярному вектору $\hat{\bm{A}}$, отвечающему минимальному сингулярному числу.

% ==============================================================================
\section{Вычислительная устойчивость и регуляризация}
\label{sec:stability}
% ==============================================================================

На практике при решении систем $\bm{A}\bm{x} = \bm{b}$ возникают проблемы, связанные с конечной машинной точностью и плохой обусловленностью матрицы системы.

\begin{tcbexample}[Матрица Гильберта]
\label{ex:hilbert}
    Элементы матрицы Гильберта $\bm{H} \in \R^{m \times n}$ задаются формулой $h_{ij} = \frac{1}{i+j-1}$.
    Сингулярные числа такой матрицы экспоненциально убывают к нулю. Несмотря на то что формально матрица невырождена (при $m=n$), младшие сингулярные векторы практически неразличимы от машинного шума. При попытке решить систему $\bm{H}\bm{x} = \bm{b}$ напрямую, относительная ошибка решения может достигать критических значений.
\end{tcbexample}

\begin{tcbinsight}[Усеченное сингулярное разложение]
    При разложении зашумленной правой части $\bm{b}$ по базису из левых сингулярных векторов $\bm{U}$ компоненты, соответствующие сверхмалым сингулярным числам $\sigma_i$, содержат преимущественно шум. 
    
    Для стабилизации решения применяется усеченное сингулярное разложение (Truncated SVD). Задавая порог отсечения $\delta > 0$ (например, $10^{-7}$), мы обращаем в нуль все $\sigma_i < \delta$ при вычислении псевдообратной матрицы $\bm{A}^+$. Искусственное зануление младших компонент подавляет усиление шума и действует как метод регуляризации Тихонова, кардинально снижая относительную ошибку полученного псевдорешения.
\end{tcbinsight}

% ==============================================================================
\section{Ортогонализация: Хаусхолдер и Грама--Шмидт}
\label{sec:orthogonalization}
% ==============================================================================

Ортогонализация векторов и вычисление QR-разложения подвержены влиянию ошибок округления. Классический алгоритм Грама--Шмидта (CGS) при применении к плохо обусловленным матрицам (см.~\cref{ex:hilbert}) катастрофически теряет ортогональность выходных векторов (норма $\norm{\bm{I} - \bm{Q}\T\bm{Q}}_2$ значительно отличается от нуля).

\vspace{5mm}
\subsection{Отражения Хаусхолдера}
\label{subsec:householder}

\begin{tcbdefinition}[Матрица Хаусхолдера]
\label{def:householder}
    Для вектора $\bm{u} \in \C^n$ такого, что $\norm{\bm{u}}_2 = 1$, матрица отражения Хаусхолдера задается выражением
    \begin{equation}
        \bm{H}(\bm{u}) = \bm{I} - 2\bm{u}\bm{u}\H.
        \label{eq:householder}
    \end{equation}
\end{tcbdefinition}
Матрица Хаусхолдера является унитарной ($\bm{H}\H \bm{H} = \bm{I}$) и эрмитовой ($\bm{H}\H = \bm{H}$). 

\begin{tcbtheorem}[Вычислительная сложность QR через Хаусхолдера]
\label{thm:householder-complexity}
    Сложность построения полного QR-разложения матрицы $\bm{A} \in \R^{n \times n}$ с использованием отражений Хаусхолдера составляет асимптотически $\frac{4}{3}n^3$ арифметических операций.
\end{tcbtheorem}

\begin{tcbproof}[теоремы~\ref{thm:householder-complexity}]
    На $k$-м шаге алгоритма происходит умножение матрицы Хаусхолдера $\bm{H}(\bm{u}_k)$ на текущую подматрицу $\tilde{\bm{A}} \in \R^{(n-k) \times (n-k)}$:
    \begin{equation}
        \bm{H}(\bm{u}_k)\tilde{\bm{A}} = \tilde{\bm{A}} - 2\bm{u}_k(\bm{u}_k\T \tilde{\bm{A}}).
        \label{eq:householder-step}
    \end{equation}
    Подсчет операций для одного шага:
    \begin{itemize}
        \item Вычисление $\bm{u}_k\T \tilde{\bm{A}}$ требует порядка $(n-k)^2$ умножений и сложений (суммарно $2(n-k)^2$).
        \item Внешнее произведение и вычитание требует еще порядка $2(n-k)^2$ операций.
    \end{itemize}
    Следовательно, сложность одного шага составляет $\approx 4(n-k)^2$ операций. Суммируя по всем шагам от $0$ до $n-1$:
    \begin{align*}
        \sum_{k=0}^{n-1} 4(n-k)^2 &= 4 \sum_{t=1}^{n} t^2 = 4 \frac{n(n+1)(2n+1)}{6} \\
        &\sim \frac{8n^3}{6} = \frac{4}{3}n^3 + O(n^2).
    \end{align*}
    Для сравнения, стандартный алгоритм Грама--Шмидта требует $\approx 2n^3$ операций.
\end{tcbproof}

\begin{tcbwarning}[Выбор знака при построении $\bm{u}$]
\label{warn:householder-sign}
    При выборе вектора $\bm{u}$ для обнуления элементов столбца $\bm{a}_1$ используется формула $\bm{v} = \bm{a}_1 \pm \norm{\bm{a}_1}_2 \bm{e}_1$.
    При реализации на ЭВМ КРИТИЧЕСКИ ВАЖНО выбирать знак, совпадающий со знаком первого элемента $\bm{a}_{11}$, то есть:
    \[
        \bm{v} = \bm{a}_1 + \sign(\bm{a}_{11}) \norm{\bm{a}_1}_2 \bm{e}_1.
    \]
    Выбор противоположного знака может привести к катастрофической потере точности (catastrophic cancellation) при вычитании близких чисел.
\end{tcbwarning}

\vspace{5mm}
\subsection{Модифицированный алгоритм Грама--Шмидта (MGS)}
\label{subsec:mgs}

Для повышения численной устойчивости классического алгоритма применяется его модифицированная версия. Идея заключается в том, что проекция вектора вычитается не на этапе формирования нового ортогонального вектора, а немедленно вычитается из \textit{всех оставшихся} необработанных векторов.

В точной арифметике подходы идентичны, однако на ЭВМ условие $\bm{q}_i\T \bm{q}_j = 0$ выполняется лишь приближенно. MGS учитывает эти малые отклонения, корректируя последующие проекции.

\begin{tcbalgorithmbox}[Модифицированный алгоритм Грама--Шмидта (MGS)]
\label{alg:mgs}
    \textbf{Вход:} Линейно независимые векторы $\bm{a}_1, \bm{a}_2, \ldots, \bm{a}_n \in \R^m$.\\
    \textbf{Выход:} Ортонормированная система векторов $\bm{q}_1, \bm{q}_2, \ldots, \bm{q}_n$.
    \begin{enumerate}
        \item $\bm{q}_1 \gets \bm{a}_1 / \norm{\bm{a}_1}_2$.
        \item Для $k = 1, 2, \ldots, n-1$ выполнить:
        \begin{enumerate}
            \item Для $j = k+1, \ldots, n$ вычесть из $\bm{a}_j$ компоненту вдоль $\bm{q}_k$:
            \[
                \bm{a}_j \gets \bm{a}_j - \inner{\bm{a}_j, \bm{q}_k}\bm{q}_k.
            \]
            \item Нормировать текущий вектор для получения следующего элемента базиса:
            \[
                \bm{q}_{k+1} \gets \bm{a}_{k+1} / \norm{\bm{a}_{k+1}}_2.
            \]
        \end{enumerate}
    \end{enumerate}
\end{tcbalgorithmbox}

